\section{Discussion}
\label{sec:discussion}

\subsection{Relation to Prior Work}

The review points to an underdeveloped research object: the preservation of
assumptions and evidence between existing planners, predictors, monitors, and
safety mechanisms.  The framework makes these transfers explicit, and it is
compatible with, but no substitute for, the mission, component, dynamical, and
runtime-assurance theories of \cref{sec:taxonomy}.  The local condition
\eqref{eq:compatibility} does not prove a global mission property.  Feedback
loops, shared variables, circular assumptions, realizability, probabilistic
dependence, lifecycle ordering, and authority transfer require a formal
semantics chosen for the implementation at hand.

Within the reviewed corpus, the coordinated-autonomy work of Chaki, Dolan,
Giampapa, Porter, and colleagues is the closest mission-level predecessor
because it explicitly relates robot coordination to quantitative assurance and
reliability
\cite{chaki2013assurance,chaki2013probabilistic,porter2013reliability,
stancliff2006missionreliability}.  The interface-contract framework works at a
different level of abstraction: it inventories the obligations that interfaces
must carry before a mission-level model, monitor, or reliability estimate can be
applied consistently to a heterogeneous deployed stack.  A natural connection
is to treat contract rejections, recoveries, and evidence lineage as observable
variables of a quantitative mission model.  RCL is the closest component-level
counterpart: it gives behavioral contracts formal semantics and generates
monitors \cite{luckcuck2026compositional}, whereas the framework supplies the
inventory of payload, validity, evidence, compatibility, admission, and recovery
obligations that such a formalism must instantiate.  PhyAgentOS and the
verification-aware swarm fabric of Kekatos \emph{et al.}\ implement parts of
this inventory inside one runtime or one messaging layer
\cite{liu2026phyagentos,kekatos2026swarmrta}; the framework states the
obligations independently of any architecture so that they can be checked on
both vertical and horizontal interfaces.

\subsection{Guidance for ROS~2 Deployments}

A ROS~2 developer can apply the framework by auditing every inter-robot topic,
service, and action, together with the vertical interfaces that feed them.  A
ROS~2 message header carries a time stamp, but whether it records acquisition,
processing, or publication time is a convention of the producer rather than a
checked property \cite{macenski2022ros2}.  The audit should therefore establish
that peer state carries its event time and frame, that its clock domain and an
offset bound are known, and that the meaning of any covariance or confidence is
declared.  The interface should identify the configuration, the middleware
implementation and version, and the team epoch.  The consumer should evaluate
an admission predicate such as \eqref{eq:temporal-admission} and invoke a
declared recovery action instead of continuing silently on stale or
incompatible input.  The failures in \cref{sec:middleware-evidence} add two
lifecycle rules: initialize and shut down middleware in an explicit order rather
than from static destructors, and treat the discovery of a service as distinct
from the ability to complete a request--response exchange.  Not every message
needs every field; the downstream component's assumptions determine which are
required.

\subsection{Limitations}

The evidence base combines formal analyses, simulations, standards, reviews,
issue reports, and robotics demonstrations created for different purposes.
There is no shared benchmark from which to estimate prevalence, effect size, or
comparative reliability.  The search was structured but not systematic, and
terminology varies across formal methods, robotics, distributed systems, and
embodied AI.  The coverage matrix exposes coverage and evidence boundaries, but
it cannot show that the framework improves mission outcomes.

The review also concentrates on the command-to-action path.  It does not treat
security, privacy, human factors, certification, or the breadth of multi-robot
planning algorithms, each of which can add contract obligations.  Security
intersects the identity, provenance, and authority fields but requires a
separate threat model.  Resource availability appears as an operational
assumption; resource-contract theory is not reviewed in depth.
\Cref{prop:temporal-admission} covers one admission predicate for one coupled
constraint; other constraints need their own derivations.

\subsection{Open Problems}

Four problems stand out.  First, the reference fields need formal semantics that
compose across discrete mission logic, continuous dynamics, probabilistic
prediction, and asynchronous communication.  Second, lifecycle contracts must
handle membership, partition recovery, ownership, and retained state without
assuming a globally consistent instant.  Third, uncertainty and evidence need
machine-checkable types so that a consumer can decide whether a proof, a
calibration result, or a deployment test supports its decision.  Fourth,
monitoring must scale without turning every peer exchange into an unbounded
evidence log or a centralized bottleneck.

Evaluation should test the framework as an executable mechanism rather than
count populated metadata fields.  Relevant outcomes include false admission and
false rejection, time to detection, correctness of fault attribution, recovery
time, task completion, unsafe events, and runtime overhead.  Reproducible studies
should report the team size, platform, software and middleware versions,
configuration, timing assumptions, injected faults, and run-level traces.  Such
work is necessary before the framework can support a reliability or safety
claim.
